\chapter*{}
\thispagestyle{empty}

\begin{center}
       {\large\bfseries \myTitle}\\
\end{center}
\begin{center}
       \myName\\
\end{center}

\vspace{0.7cm}
\noindent{\textbf{Palabras clave}: metaheurísticas, modelos subrogados,\\
       aprendizaje automático, hibridación, optimización continua, optimización costosa.}\\

\vspace{0.7cm}
\noindent{\textbf{Resumen}}\\

La optimización mediante metaheurísticas permite abordar problemas complejos
para los que los métodos exactos no resultan viables. Sin embargo, estos
algoritmos suelen requerir numerosas evaluaciones de la función objetivo.
Cuando estas implican ejecutar simulaciones o procesos costosos, su
aplicabilidad práctica queda limitada.

Los modelos subrogados permiten reducir este coste al aproximar la función
objetivo a partir de las soluciones evaluadas durante la búsqueda. Su integración no es inmediata, ya que la distribución de estas soluciones cambia conforme evoluciona la población y puede verse condicionada por fases de estancamiento. Antes de incorporar un modelo al bucle de optimización, resulta necesario analizar qué datos generan las metaheurísticas y cuándo conservan información útil para orientar la búsqueda.

Este trabajo estudia el uso de modelos de aprendizaje automático como técnicas
subrogadas para su integración en metaheurísticas evolutivas. Se analizan las
trayectorias generadas por AGE, DE y SHADE sobre CEC2017 y se evalúa un
mecanismo de reinicio elitista destinado a reactivar la exploración. A partir
de las muestras obtenidas, se comparan diferentes familias de modelos mediante
estrategias \textit{offline} que respetan la estructura temporal de la búsqueda.
El estudio considera su capacidad para preservar el orden relativo entre
soluciones y su coste computacional, con el fin de identificar las
configuraciones más adecuadas para una posterior integración \textit{online}.

\cleardoublepage

\thispagestyle{empty}

\begin{center}
       {\large\bfseries Use of Machine Learning Techniques for the Hybridization of Metaheuristics}\\
\end{center}
\begin{center}
       \myName\\
\end{center}

\vspace{0.7cm}
\noindent{\textbf{Keywords}: metaheuristics, surrogate models,\\
       machine learning, hybridization, continuous optimization, expensive optimization.}\\

\vspace{0.7cm}
\noindent{\textbf{Abstract}}\\

Optimization through metaheuristics makes it possible to address complex
problems for which exact methods are not feasible. However, these algorithms
usually require numerous objective function evaluations. When these involve
running simulations or costly processes, their practical applicability
becomes limited.

Surrogate models can reduce this cost by approximating the objective function
from the solutions evaluated during the search. Their integration is not straightforward, as the distribution of these solutions changes as the population evolves and may be affected by stagnation phases. Before incorporating a model
into the optimization loop, it is therefore necessary to analyse the data
generated by the metaheuristics and determine when they retain useful
information for guiding the search.

This work studies the use of machine learning models as surrogate techniques
for their integration into evolutionary metaheuristics. The trajectories
generated by AGE, DE and SHADE on CEC2017 are analysed, and an elitist restart
mechanism designed to reactivate exploration is evaluated. Based on the
resulting samples, different families of models are compared using
\textit{offline} strategies that preserve the temporal structure of the search
process. The study considers their ability to preserve the relative ordering
of solutions and their computational cost in order to identify the most
suitable configurations for subsequent \textit{online} integration.

\chapter*{}
\thispagestyle{empty}

\noindent\rule[-1ex]{\textwidth}{2pt}\\[4.5ex]

Yo, \textbf{\myName}, alumno de la titulación \myDegree de la \textbf{\myFaculty de la \myUni}, con DNI 75941131F, autorizo la ubicación de la siguiente copia de mi Trabajo Fin de Grado en la biblioteca del centro para que pueda ser consultada por las personas que lo deseen.

\vspace{6cm}

\noindent Fdo: \myName

\vspace{2cm}

\begin{flushright}
       Granada, a 22 de junio de 2026.
\end{flushright}

\chapter*{}
\thispagestyle{empty}

\noindent\rule[-1ex]{\textwidth}{2pt}\\[4.5ex]

D. \textbf{\myProf}, Profesor del \textbf{\myDepartment} de la \textbf{\myUni}.

\vspace{0.5cm}

\textbf{Informa:}

\vspace{0.5cm}

Que el presente trabajo, titulado \textit{\textbf{\myTitle}}, ha sido realizado bajo su supervisión por \textbf{\myName}, y autoriza la defensa de dicho trabajo ante el tribunal que corresponda.

\vspace{0.5cm}

Y para que conste, expide y firma el presente informe en Granada, 22 de junio de 2026.

\vspace{1cm}

\textbf{El director:}

\vspace{5cm}

\noindent \textbf{\myProf}

\chapter*{Agradecimientos}
\thispagestyle{empty}

\vspace{1cm}

% En primer lugar, quiero agradecer y dedicar este proyecto a mis padres, ya que gracias a ellos he podido aprender y formarme en aquello que siempre me había gustado. Gracias por apoyarme en todo momento cuando lo he necesitado, por educarme para ser una persona trabajadora y esforzada, capaz de dar siempre lo máximo para conseguir todo aquello que se propone, y por las incontables horas de llamada a lo largo de estos cuatro años, mientras me quejaba, me quejaba y me quejaba.

% Gracias también a mi pareja, presente en todo momento, tanto en los buenos como en los malos, incluso en aquellos en los que el esfuerzo no llega a verse recompensado como uno esperaba, y capaz de animarme con muy pocas palabras.

% También quiero agradecer a todos los amigos que he hecho en la universidad, que han sido y seguirán siendo un pilar muy importante durante estos años. Gracias por todo el apoyo y la ayuda que siempre nos hemos querido dar sin esperar nada a cambio, por todas las risas y, por supuesto, por haber conseguido que una relación de compañeros de carrera llegase a convertirse en una gran amistad. Pero, sobre todo, quiero agradecer a Francesc Oliver Catany, Juan Carlos Mora Díaz, Jorge López Molina y Mehdi Iabouten que, a pesar de la distancia, nunca ha perdido el interés por nuestra relación.
